%!TeX root=MemoriaTFG.tex

\chapter*{Resumen}
\addcontentsline{toc}{chapter}{Resumen}
\markboth{RESUMEN}{RESUMEN}
\vspace*{0.4cm}

La dermatología clínica afronta una demanda asistencial creciente frente a un acceso limitado al especialista, lo que ha impulsado el interés por los sistemas de apoyo al diagnóstico basados en imagen. En los últimos años han aparecido \emph{modelos fundacionales} ---modelos de gran escala, preentrenados y reutilizables mediante adaptación--- específicos para dermatología; sin embargo, su rendimiento real fuera de los conjuntos con los que se desarrollaron, y muy especialmente sobre material clínico en español, apenas se ha evaluado de forma independiente.
\par Este trabajo aporta tres contribuciones cerradas y reproducibles. Primero, una \textbf{evaluación comparativa homogénea} de los principales modelos fundacionales dermatológicos con pesos públicos, sobre un conjunto común de tareas ---clasificación, segmentación, recuperación de casos y clasificación \emph{zero-shot} guiada por texto--- y de datasets públicos armonizados mediante una ontología clínica de tres niveles. Segundo, \textbf{DermapixelAI 1.0}, un dataset clínico en castellano con texto narrativo asociado a cada caso, entregado con documentación formal. Tercero, \textbf{Dermapixel R0}, una adaptación propia que demuestra que un modelo fundacional puede especializarse al castellano con recursos modestos.
\par Leídos en conjunto, los resultados no señalan un modelo universalmente superior, sino una especialización por tarea: la elección depende del tipo de problema, del volumen de datos etiquetados y de la disparidad entre fototipos cutáneos. De aquí se desprende una idea central: una vez los codificadores visuales alcanzan cierta madurez, el principal factor limitante deja de ser la arquitectura y pasa a ser la calidad, diversidad y cobertura de los datos disponibles. El escalado del dataset con material hospitalario y la validación clínica prospectiva constituyen la continuación natural del trabajo para su futura integración en sistemas de teledermatología hospitalaria.

\vspace*{-0.4cm}
\section*{Palabras clave}
modelos fundacionales, dermatología clínica, PanDerm, DermLIP,
evaluación empírica, ontología, equidad, dataset en español.

% ----------------------------------------------------------------------

\chapter*{Abstract}
\addcontentsline{toc}{chapter}{Abstract}
\markboth{ABSTRACT}{ABSTRACT}
\vspace*{0.4cm}

Clinical dermatology faces growing demand against limited access to specialists, which has driven interest in image-based diagnostic support systems. Recent years have seen the emergence of \emph{foundation models} ---large, pretrained models reusable through adaptation--- specific to dermatology; their real-world performance beyond the datasets they were developed on, and particularly on Spanish-language clinical material, has barely been evaluated independently.
\par This work delivers three closed, reproducible contributions. First, a \textbf{homogeneous comparative evaluation} of the leading dermatology foundation models with public weights, across a common set of tasks ---classification, segmentation, case retrieval and text-guided zero-shot classification--- and public datasets harmonised through a three-level clinical ontology. Second, \textbf{DermapixelAI 1.0}, a Spanish clinical dataset with narrative text linked to each case, released with formal documentation. Third, \textbf{Dermapixel R0}, an in-house adaptation showing that a foundation model can be specialised to Spanish with modest resources.
\par Taken together, the results point not to a universally superior model but to a specialisation by task: the right choice depends on the problem type, the labelled-data budget and skin-phototype disparity. A central idea follows: once visual encoders reach a certain maturity, the main limiting factor shifts from architecture to the quality, diversity and coverage of the available data. Scaling the dataset with hospital material and prospective clinical validation constitute the natural continuation of this work towards its future integration into hospital teledermatology systems.

\vspace*{-0.4cm}
\section*{Keywords}
foundation models, clinical dermatology, PanDerm, DermLIP, empirical
evaluation, ontology, fairness, Spanish dataset.
